\documentclass[10pt,a4paper]{amsart}
\usepackage[margin=19mm]{geometry}
\usepackage{amsmath,amssymb,mathtools}
\usepackage[scr=rsfs,cal=euler]{mathalfa}
\usepackage{xfrac}
\usepackage[unicode,hidelinks,bookmarks=false]{hyperref}
\newcommand{\ie}{i.e.\ }
\newcommand{\cf}{cf.\ }
\newcommand{\resp}{respectively\ }
\newcommand{\Subsection}{\subsection}
\providecommand{\nobreakdash}{\nobreak\hbox{-}}
\setlength{\emergencystretch}{2em}
\multlinegap0pt
\usepackage{amssymb}
\usepackage[shortlabels]{enumitem}
\newtheorem{theorem}{Theorem}
\title{Finite local metric models for residually finite isometric actions}
\author{Vadim Alekseev, Andreas Thom}
\date{Source: arXiv:2512.14147v1 (CC BY 4.0)}
\begin{document}
\maketitle

\noindent\emph{Source and licence.} On finite local approximations of isometric actions of residually finite groups. arXiv:2512.14147v1 (CC BY 4.0), https://arxiv.org/abs/2512.14147v1. The exact-version arXiv abstract page was independently checked on 2026-10-10 and declares CC BY 4.0. \texttt{licenses/arxiv-2512.14147-CC-BY-4.0.txt}.\par\smallskip

\noindent\textit{Editorial scope.} The counted unit is Theorem thm:main, source0.tex lines 66--72, with the entire author proof at lines 96--248 copied without mathematical alteration. Lines 92--94 retain the definition of residual finiteness and the authors\textquotesingle{} attribution of proof inspiration; the native bibliography entry at lines 305--308 is retained. The title of this delivery describes the selected theorem; the source work is named above. The authors\textquotesingle{} manuscript-level AI-assistance disclosure at line 281 is reproduced after the proof. Under the explicitly approved collection policy, authors remain responsible for the final substantially revised proof; this material is not labelled AI-free. No extra proof is generated. The unchanged native source is bundled as \texttt{sources/191078/source0.tex}.
The following theorem is the main result. Recall, a group $G$ is \emph{residually finite} if for every $g\in G\setminus\{e\}$ there exists a finite group $Q$ and a homomorphism $\psi\colon G\to Q$ with $\psi(g)\neq e$. The proof is inspired
by Doucha's work on metric approximation of topological groups,
see in particular~\cite[Proposition 3.9]{DouchaMetricTopGroups}.

\begin{theorem} \label{thm:main}
Let $G$ be a residually finite group and let $G \curvearrowright_{\varphi} (X,d)$ be an isometric action. For all $X_0 \subset X$ finite, $A \subset G$ finite and $\varepsilon>0$, there exists a finite metric space $(Y,\eta)$ with an isometric action $G \curvearrowright_\psi (Y,\eta)$ and a map $f \colon X_0 \to Y$, such that
\[
\bigl|d(\varphi_g(x),\varphi_h(y)) - \eta(\psi_g(f(x)),\psi_h(f(y)))\bigr|\leq \varepsilon,
\quad \forall g,h \in A, \ x,y\in X_0.
\]
\end{theorem}

\begin{proof}[Proof of Theorem \ref{thm:main}]
Without loss of generality we may assume that $A=A^{-1}$ and $e \in A$.
Let $X_0\subset X$ be finite and set $Z:=G \times X_0.$
Define a pseudometric $\kappa$ on $Z$ by
$\kappa((g,x),(h,y)) := d(\varphi_g(x),\varphi_h(y)).$
Because the action is isometric, $\kappa$ is left $G$-invariant in the first coordinate, i.e.,
$\kappa((kg,x),(kh,y)) = \kappa((g,x),(h,y))$ for all $k,g,h\in G,\ x,y\in X_0$.
Let $d_{\rm st}\colon Z\times Z\to\{0,1\}$ be the discrete metric
\[
d_{\rm st}((g,x),(h,y)) =
\begin{cases}
0 &\text{if } (g,x)=(h,y),\\
1 &\text{otherwise.}
\end{cases}
\]
Define
\[
d_{\varepsilon}((g,x),(h,y)):=\kappa((g,x),(h,y)) + \varepsilon d_{\rm st}((g,x),(h,y)).
\]
Then $d_\varepsilon$ is a genuine metric on $Z$, still left-invariant in the first coordinate. In particular, for $(g,x)\neq(h,y)$ we have $d_\varepsilon((g,x),(h,y))\ge\varepsilon$.
Set $m:= \max \{d_{\varepsilon}((g,x),(h,y)) \mid x,y \in X_0,\ g,h \in A\}$
and choose $k\in \mathbb N$ with $k \geq 2\varepsilon^{-1}m.$
Now, using that $G$ is residually finite, fix a finite quotient
$ \pi\colon G\to H$
such that $\pi$ is injective on the finite set $B:=A^{k}$. Let $H_0:=\langle \pi(A)\rangle\le H$ be the subgroup generated by $\pi(A)$.

We set $Y_0:=H_0\times X_0$ and define a weighted graph structure on $Y_0$ as follows: we declare that there is an edge between the vertices
$(q\pi(g),x)$ and $(q\pi(h),y)$, for $q\in H_0$, $g,h\in A$, $x,y\in X_0$, with weight
$w\bigl((q\pi(g),x),(q\pi(h),y)\bigr)
:= d_\varepsilon((g,x),(h,y)).$

Let $\eta$ be the induced path metric on $Y_0$: for any two vertices, $\eta$ is the infimum of the sums of edge weights over all finite paths joining them. By construction, left multiplication by an element of $H_0$ on the first coordinate permutes vertices and preserves edge weights, so the action
$H_0\curvearrowright Y_0$, $h\cdot(q,x):=(hq,x)$,
is by isometries. We can extend this metric to $Y:=H \times X_0$ in a $H$-invariant way that sets all distances which are not fixed by $H$-invariance equal to a large constant. This way, 
we obtain an isometric action of  $G$ on $(Y,\eta)$ by
$\psi_g(q,x):=(\pi(g)q,x).$ Finally, define $f\colon X_0\to Y$ by $f(x):=(e,x)$ where $e\in H$ is the identity.

We need to compare $\eta$ and $d_\varepsilon$ on $A\times X_0$. We now show that for all $g,h\in A$, $x,y\in X_0$,
\begin{equation}\label{eq:eta-deps}
\eta\bigl(\psi_g(f(x)),\psi_h(f(y))\bigr)
= d_\varepsilon((g,x),(h,y)).
\end{equation}
This amounts to proving
\[
\eta\bigl((\pi(g),x),(\pi(h),y)\bigr) = d_\varepsilon((g,x),(h,y)).
\]

For $g,h\in A$ and $x,y\in X_0$ there is an edge between $(\pi(g),x)$ and $(\pi(h),y)$ corresponding to $q=e$, so
\[
\eta\bigl((\pi(g),x),(\pi(h),y)\bigr)
\le w\bigl((\pi(g),x),(\pi(h),y)\bigr)
= d_\varepsilon((g,x),(h,y)).
\]
We prove the reverse inequality by contradiction. Fix $g,h\in A$, $x,y\in X_0$ and suppose that
\begin{equation}\label{eq:contr}
\eta\bigl((\pi(g),x),(\pi(h),y)\bigr)
< d_\varepsilon((g,x),(h,y)).
\end{equation}
Let
\[
(\pi(g),x)=(v_0,x_0), (v_1,x_1),\dots,(v_n,x_n)=(\pi(h),y)
\]
be a path in $Y$ with minimal number $n$ of edges and such that
\[
\eta\bigl((\pi(g),x),(\pi(h),y)\bigr)
= \sum_{j=1}^n w_j,
\]
where $w_j$ is the weight of the $j$-th edge. Each edge comes from some triple $q_j\in H_0,g_j,h_j\in A$,
and has the form
\[
(v_{j-1},x_{j-1})=(q_j\pi(g_j),x_{j-1}),\qquad
(v_j,x_j)=(q_j\pi(h_j),x_j),
\]
with $w_j = d_\varepsilon((g_j,x_{j-1}),(h_j,x_j))$.
Since no edge is trivial, we have $w_j\ge\varepsilon$, hence
\[
\varepsilon n \le \sum_{j=1}^n w_j
= \eta\bigl((\pi(g),x),(\pi(h),y)\bigr)
< d_\varepsilon((g,x),(h,y))
\le m,
\]
and therefore $n < \varepsilon^{-1}m \le k/2$.
Next, consider the elements
\[
u := g^{-1}h,\qquad
u' := \prod_{j=1}^n g_j^{-1}h_j \in G.
\]
We claim that
\begin{equation}\label{eq:uuprime}
\pi(u) = \pi(u').
\end{equation}
Indeed, from $v_{j-1}=q_j\pi(g_j)$ and $v_j=q_j\pi(h_j)$ we get
$v_{j-1}^{-1}v_j = \pi(g_j^{-1}h_j),$
and hence
\[
v_n
= v_0 \prod_{j=1}^n v_{j-1}^{-1}v_j
= \pi(g)\,\prod_{j=1}^n \pi(g_j^{-1}h_j)
= \pi\Bigl(g \prod_{j=1}^n g_j^{-1}h_j\Bigr)
= \pi(gu').
\]
Since $v_n=\pi(h)$, this gives $\pi(gu')=\pi(h)$, i.e.~$\pi(u')=\pi(g^{-1}h)=\pi(u)$.

Now observe that $u = g^{-1}h \in A^2 \subset A^{k}=B$,
and each $g_j^{-1}h_j\in A^2$, so $u'=\prod_{j=1}^n g_j^{-1}h_j$ lies in $A^{2n}\subset A^{k}=B$, because $n\le k/2$. Thus $u,u'\in B$ and, by injectivity of $\pi$ on $B$, \eqref{eq:uuprime} implies $u'=u$, i.e.
\begin{equation}\label{eq:gprod= h}
g\prod_{j=1}^n g_j^{-1}h_j = h.
\end{equation}

Define elements $b_0,\dots,b_n\in G$ recursively by
$b_0:=g, b_j := b_{j-1}g_j^{-1}h_j\quad (1\le j\le n).$ Then by construction
\[
b_n = g\prod_{j=1}^n g_j^{-1}h_j = h
\]
by \eqref{eq:gprod= h}. We now compare $w_j$ with $d_\varepsilon$ along the sequence $(b_j,x_j)$. For each $j$ we have
$b_{j-1}^{-1}b_j = g_j^{-1}h_j.$ Using left-invariance of $d_\varepsilon$ we obtain
\begin{align*}
d_\varepsilon((b_{j-1},x_{j-1}),(b_j,x_j))
&= d_\varepsilon\bigl((e,x_{j-1}),(b_{j-1}^{-1}b_j,x_j)\bigr)\\
&= d_\varepsilon\bigl((e,x_{j-1}),(g_j^{-1}h_j,x_j)\bigr)\\
&= d_\varepsilon((g_j,x_{j-1}),(h_j,x_j))\\
&= w_j.
\end{align*}
Therefore
\[
\sum_{j=1}^n w_j
= \sum_{j=1}^n d_\varepsilon((b_{j-1},x_{j-1}),(b_j,x_j)).
\]
By the triangle inequality for $d_\varepsilon$ and the facts that $(b_0,x_0)=(g,x)$ and $(b_n,x_n)=(h,y)$, we get
\[
d_\varepsilon((g,x),(h,y))
\le \sum_{j=1}^n d_\varepsilon((b_{j-1},x_{j-1}),(b_j,x_j))
= \sum_{j=1}^n w_j
= \eta\bigl((\pi(g),x),(\pi(h),y)\bigr),
\]
which contradicts our assumption \eqref{eq:contr}. Hence no such path exists, and we must have
\[
\eta\bigl((\pi(g),x),(\pi(h),y)\bigr)
= d_\varepsilon((g,x),(h,y)).
\]
Finally, for $g,h\in A$ and $x,y\in X_0$ we have
\[
\eta(\psi_g(f(x)),\psi_h(f(y)))
= d_\varepsilon((g,x),(h,y))
= d(\varphi_g(x),\varphi_h(y)) + \varepsilon\, d_{\rm st}((g,x),(h,y)),
\]
and $d_{\rm st}\in\{0,1\}$, so
\[
\bigl|\eta(\psi_g(f(x)),\psi_h(f(y))) - d(\varphi_g(x),\varphi_h(y))\bigr|
\le \varepsilon.
\]
This completes the proof.
\end{proof}


\noindent\textit{Author disclosure from the source article.}
The second author thanks the Isaac Newton Institute for its hospitality during the workshop \emph{Stability and probabilistic methods} in November 2025. ChatGPT was used to assist in drafting parts of this manuscript. All content was reviewed and substantially revised by the authors, who are responsible for the final text. We thank Michal Doucha for interesting remarks on a first draft of this paper.

\begin{thebibliography}{99}
\bibitem[4]{DouchaMetricTopGroups} M.~Doucha, \emph{Metric topological groups: their metric approximation and metric ultraproducts}, Groups Geom. Dyn. \textbf{12} (2018), no.~2, 615--636.
\end{thebibliography}
\end{document}
